\chapter{Posicionamento Exploratório contra Outros Comparadores}
\label{apend:exploratorio}

Este apêndice registra o posicionamento do IVF/SPEA2 contra os sete comparadores de contexto. Ele fica fora dos capítulos de resultados por duas razões: o posicionamento é descritivo, sem teste de hipótese nem família de hipóteses declarada no protocolo, e a pergunta que responde (qual algoritmo ordena melhor no conjunto avaliado) não é a QP1, respondida na Seção~\ref{sec:confirmatorio}. Nenhuma questão de pesquisa desta dissertação depende deste resultado, que não substitui nem reforça a comparação principal.

Uma análise de posto médio entre os nove algoritmos coloca o IVF/SPEA2 em segundo lugar com $M = 2$ e em primeiro com $M = 3$. Em cada instância, os algoritmos são ordenados pela mediana de IGD, e o posto médio agrega essas ordens sobre instâncias heterogêneas, sem considerar a magnitude das diferenças nem a variabilidade entre execuções; o resultado descreve o conjunto avaliado e não sustenta afirmação de superioridade geral. A Figura~\ref{fig:posto_medio} apresenta esse posto médio, com o desvio padrão entre instâncias, para $M = 2$ e $M = 3$.

\begin{figure}[!htb]
    \centering
    \includegraphics[width=\textwidth]{generated/figures/fig_posto_medio.pdf}
    \caption[Posto médio exploratório de IGD entre os nove algoritmos]{Posto médio exploratório de IGD entre os nove algoritmos avaliados, calculado sobre os postos das medianas em cada instância (28 instâncias em $M = 2$; 23 em $M = 3$), mesma coorte de submissão da Seção~\ref{sec:estatistica}, $n=60$ execuções por algoritmo e instância. IGD: menor é melhor, logo menor posto é melhor; barras horizontais são o desvio padrão do posto entre instâncias. Descritivo, sem teste de hipótese.}
    \label{fig:posto_medio}
\end{figure}
